% =====================================================================
% TikZ 模板骨架（步骤③加载）
%
% 用法：复制本文件为新 .tex，按需修改。所有 ⚠️ 注释为防御性写法，禁删。
%
% ⚠️ 编译前必须 `fc-list | grep "PingFang SC"` 确认 CJK 字体存在
%    平台优先级：macOS → PingFang SC; Linux → Noto Sans CJK SC;
%               Windows → SimHei / Microsoft YaHei
% ⚠️ 编译后必须 `grep "Missing character" *.log`——xelatex 对缺失字体静默失败
% =====================================================================

\documentclass[tikz,border=25pt]{standalone}
\usepackage{tikz}
% Overleaf 编译时可替换为 \usepackage{ctex}
% 方案 B（无 ctex 时）：\usepackage{fontspec} + \setmainfont{...} + \setsansfont{...}
% ⚠️ 不要用 ucharclasses 方案——中英混排场景下 Missing character 高发
\usepackage[fontset=none]{ctex}
\setCJKmainfont{PingFang SC}   % ← 按本机可用字体替换
\setCJKsansfont{PingFang SC}   % ← 同上
\usetikzlibrary{shapes, arrows.meta, bending, positioning, fit, backgrounds, calc, shadows}
% ⚠️ bending library 必须加载 — 否则箭头 tip 在曲线/弯折路径上对齐失效（深度调研发现的根因）

% 含数学公式时额外引入：
% \usepackage{amsmath, amssymb}
% 密码学/形式化领域含 ⟦⟧ 时引入：
% \usepackage{stmaryrd}

% =====================================================================
% 配色（学术配色，默认）—— 详见 references/tikz-colors.md
% 如需 draw.io 经典配色或仅引入子集，从 tikz-colors.md 拷贝对应块即可
% =====================================================================
\definecolor{acaBlueLine}{HTML}{6080B0}      \definecolor{acaBlueFill}{HTML}{DBEAFE}
\definecolor{acaGreenLine}{HTML}{30A060}     \definecolor{acaGreenFill}{HTML}{A0D0A0}
\definecolor{acaOrangeLine}{HTML}{D06020}    \definecolor{acaOrangeFill}{HTML}{FFE6CC}
\definecolor{acaPurpleLine}{HTML}{6020D0}    \definecolor{acaPurpleFill}{HTML}{E1D5E7}
\definecolor{acaRedLine}{HTML}{B05050}       \definecolor{acaRedFill}{HTML}{F8CECC}
\definecolor{acaGreyLine}{HTML}{666666}      \definecolor{acaGreyFill}{HTML}{F5F5F5}
\definecolor{acaGoldLine}{HTML}{D09000}      \definecolor{acaGoldFill}{HTML}{F8F8E0}
\definecolor{acaTealLine}{HTML}{009060}      \definecolor{acaTealFill}{HTML}{D1FAE5}
\definecolor{acaCyanLine}{HTML}{00B0D0}      \definecolor{acaCyanFill}{HTML}{E0F7FA}
\definecolor{acaPinkLine}{HTML}{D06090}      \definecolor{acaPinkFill}{HTML}{FCE4EC}
\definecolor{acaYellowLine}{HTML}{E0C060}    \definecolor{acaYellowFill}{HTML}{FFF8E1}
\definecolor{acaLimeLine}{HTML}{80B060}      \definecolor{acaLimeFill}{HTML}{E8F5E9}

% draw.io 兼容别名（指向学术配色）—— 旧代码可直接复用
\colorlet{drawBlueFill}{acaBlueFill}    \colorlet{drawBlueLine}{acaBlueLine}
\colorlet{drawGreenFill}{acaGreenFill}  \colorlet{drawGreenLine}{acaGreenLine}
\colorlet{drawOrangeFill}{acaOrangeFill}\colorlet{drawOrangeLine}{acaOrangeLine}
\colorlet{drawPurpleFill}{acaPurpleFill}\colorlet{drawPurpleLine}{acaPurpleLine}
\colorlet{drawRedFill}{acaRedFill}      \colorlet{drawRedLine}{acaRedLine}
\colorlet{drawGreyFill}{acaGreyFill}    \colorlet{drawGreyLine}{acaGreyLine}

% zone 背景色（极浅）
\definecolor{zoneBlueBg}{HTML}{E0E0F8}
\definecolor{zoneGreenBg}{HTML}{ECFDF5}
\definecolor{zonePurpleBg}{HTML}{F5F3FF}
\definecolor{zoneRedBg}{HTML}{F8E0E0}
\definecolor{zoneYellowBg}{HTML}{F8F0C0}
\definecolor{zoneOrangeBg}{HTML}{FFF5EB}

\pgfdeclarelayer{background}
\pgfsetlayers{background,main}

\begin{document}
\begin{tikzpicture}[
    node distance=1.2cm and 2cm,
    every node/.style={font=\footnotesize},
    base_box/.style={rectangle, rounded corners=3pt, align=center,
        minimum height=0.9cm, minimum width=2.8cm,
        inner sep=10pt,                          % 防御：文字不贴框边
        drop shadow={opacity=0.15}, thick},
    blue_node/.style={base_box, fill=acaBlueFill, draw=acaBlueLine},
    green_node/.style={base_box, fill=acaGreenFill, draw=acaGreenLine},
    orange_node/.style={base_box, fill=acaOrangeFill, draw=acaOrangeLine},
    purple_node/.style={base_box, fill=acaPurpleFill, draw=acaPurpleLine},
    red_node/.style={base_box, fill=acaRedFill, draw=acaRedLine, font=\footnotesize\bfseries},
    grey_node/.style={base_box, fill=acaGreyFill, draw=acaGreyLine},
    % 扩展配色
    gold_node/.style={base_box, fill=acaGoldFill, draw=acaGoldLine},
    teal_node/.style={base_box, fill=acaTealFill, draw=acaTealLine},
    cyan_node/.style={base_box, fill=acaCyanFill, draw=acaCyanLine},
    pink_node/.style={base_box, fill=acaPinkFill, draw=acaPinkLine},
    yellow_node/.style={base_box, fill=acaYellowFill, draw=acaYellowLine},
    lime_node/.style={base_box, fill=acaLimeFill, draw=acaLimeLine},
    % ╔════════════════════════════════════════════════════════════════╗
    % ║ 箭头/连线 canonical styles（深度调研后的最佳实践，2026-05-18）║
    % ╠════════════════════════════════════════════════════════════════╣
    % ║ 关键设计原则（来自 arrows.meta + bending library 官方推荐）：    ║
    % ║                                                                 ║
    % ║ 1. tip 大小 = `<dim> <line_width_factor>` 语法，**自动跟随线宽** ║
    % ║    例: `length=5pt 1.5` = 5pt + 1.5×line_width                  ║
    % ║    不要再用 `scale=0.75-1.3` 这种绝对 scale（4 轮 Batch 教训）   ║
    % ║                                                                 ║
    % ║ 2. width' (带 prime) = tip 宽相对 tip 长 的比例                ║
    % ║    例: `width'=0pt 0.6` = 宽是长的 60%（黄金三角比例）          ║
    % ║                                                                 ║
    % ║ 3. shorten = `<dim> <line_width_factor>`，自动跟随线宽           ║
    % ║    例: `shorten >=0pt 0.5` = 0.5×line_width                     ║
    % ║                                                                 ║
    % ║ 4. line width 是核心调节量，不是 tip scale                     ║
    % ║    1.0pt = 默认 / 1.6pt = 强调 / 0.6pt = 细节                  ║
    % ║                                                                 ║
    % ║ 5. PlotNeuralNet（业界标杆）用 `ultra thick` + 默认 Stealth     ║
    % ║    粗线+默认 tip 比 细线+缩放 tip 视觉更稳                      ║
    % ╚════════════════════════════════════════════════════════════════╝
    arrow/.style={
        -{Stealth[length=5pt 1.5, width'=0pt 0.6, sep=0pt 0.5]},
        line width=1.0pt,
        shorten >=2pt, shorten <=1pt,            % ⚠️ shorten 不支持 dual-arg；
                                                 % 用绝对 pt（TeX Live 2024 实测验证）
        color=black!70,
    },
    arrow thick/.style={arrow, line width=1.6pt},    % 主流 / 强调
    arrow thin/.style={arrow, line width=0.6pt},     % 细节 / 次要
    residual/.style={                                % 残差/skip 虚线
        -{Stealth[length=4pt 1.2, width'=0pt 0.6]},
        line width=0.9pt,
        densely dashed,
        rounded corners=5pt,
        shorten >=2pt,
        color=acaPurpleLine,
    },
    leader/.style={                                  % 标注引线
        densely dotted,
        line width=0.5pt,
        color=black!50,
        shorten >=2pt, shorten <=2pt,
    },
    tag/.style={font=\scriptsize, fill=white, inner sep=2pt, rounded corners=1pt},
    annot/.style={font=\footnotesize, inner sep=2pt},
    zone/.style={dashed, thick, inner sep=15pt, rounded corners=8pt},
]
% ===== 节点、连线、分区 =====

% 防御性写法（必须遵守）：
%
% 1. 箭头连接优先用 -| 或 |- （自动 L 形），不手动算中间点：
%    ✅ \draw[arrow] (A.east) -| (B.north);     % TikZ 自动算拐点
%    ✅ \draw[arrow] (A.south) |- (B.west);     % 先下再右
%    ❌ \draw[arrow] (A.east) -- (3.5,-8) -- (B.west);  % 手动拐点容易撞框边
%
% 2. 树状分叉必须用连续路径画主干+横杆，避免断连：
%    % 主干+横杆一笔画完（无 arrow、无 shorten）
%    \draw[thick, color=black!70]
%        (A.south) -- ++(0,-1.5)
%        -| ++(3,0) coordinate(right_end)
%        (A.south) ++(0,-1.5) -| ++(-3,0) coordinate(left_end);
%    % 各分支独立画（带 arrow + shorten）
%    \draw[arrow] (left_end) -- (B.north);
%    \draw[arrow] (A.south |- left_end) -- (C.north);
%    \draw[arrow] (right_end) -- (D.north);
%    ❌ 不要用多个独立 \draw 画主干+横杆——shorten 会让交叉点出现间隙
%    ❌ 不要在分叉点用 circle/rectangle 小节点——会出现"空心方块"
%
% 3. zone 标题用 label= 避免标题和内容重叠：
%    \node[zone, label={[font=\small\bfseries, yshift=3pt]above:标题}] ...
%
% 4. 装饰性 fill（网格、热力图色块、背景色）放 background 层，文字放 main 层：
%    \begin{pgfonlayer}{background}
%      \fill[blue!20] (0,0) rectangle (2,2);  % 色块在底层
%    \end{pgfonlayer}
%    \node at (1,1) {标签};                    % 文字在上层，不被遮挡

\end{tikzpicture}
\end{document}


% =====================================================================
% 可选扩展（按需启用）
% =====================================================================

% --- 几何/数学示意图额外需要：amsmath + 绝对坐标 + 网格 + vec 样式 ---
% \usepackage{amsmath, amssymb}
% vec/.style={->, thick, >=Stealth},
% formula_box/.style={rectangle, draw=acaGreyLine, fill=acaYellowFill!50,
%     inner sep=8pt, rounded corners=4pt, align=center}

% --- 计算机/密码学领域扩展 styles（按需添加到 tikzpicture 选项中）---
% code_block/.style={rectangle, rounded corners=2pt, draw=black!20, fill=black!3,
%     align=left, inner sep=6pt, font=\ttfamily\scriptsize, text width=4.5cm},
% diamond_node/.style={diamond, draw=drawGreyLine, fill=white, thick,
%     minimum size=1.2cm, inner sep=1pt, align=center, font=\scriptsize\bfseries},
% sum_circle/.style={circle, draw=drawGreyLine, fill=white, thick,
%     minimum size=1.2cm, font=\Large\bfseries},
% mem_cell/.style={rectangle, draw=drawGreyLine!60, fill=drawGreyFill!50,
%     minimum height=0.55cm, minimum width=2.0cm, align=center, font=\ttfamily\scriptsize},

% --- 3D 伪立体效果（仅矩形节点可用，梯形/圆/菱形会错位）---
% 在节点定义后调用，给节点添加3D效果。如需区分 Backbone 和 Head 等，用不同尺寸
% 的矩形（Head 更窄更小）而非梯形。
%
% 右侧面：
% \fill[mycolor!40, draw=mycolor!80!black, thick]
%     ([xshift=4pt,yshift=4pt]node.north east) -- (node.north east)
%     -- (node.south east) -- ([xshift=4pt,yshift=4pt]node.south east) -- cycle;
% 顶面：
% \fill[mycolor!30, draw=mycolor!80!black, thick]
%     ([xshift=4pt,yshift=4pt]node.north west) -- ([xshift=4pt,yshift=4pt]node.north east)
%     -- (node.north east) -- (node.north west) -- cycle;
